\chapter{入力はいくつ必要か}
% full version: chapters/19_dendrites.tex

ニューロンには枝分かれした入力用の突起、樹状突起がある。あるニューロンでは数千の接点をもつ1本の木まるごとであり、別のニューロンでは1、2本の枝か、まったくない。エンジニアにとって入力の数は回路のパラメータであり、課題に合わせて選ばれている。

\section*{ゼロから10万まで}

触覚と痛みのセンサには樹状突起がまったくない。配線が皮膚から脊髄へただ伸びていて、センサはその先端にある。遠端にセンサのついた線路である。

\pencilfig{19}{\pencilart{19}}{入力が多ければ加算器。巨大な入力が1つなら、速く正確な中継器。}

嗅覚のニューロンと、目や耳のセンサから信号を伝えるニューロンには、樹状突起が1本ある。入力1つ、出力1つ。リピータ、バッファである。

音の方向の検出器には、反対方向を向いた2本の樹状突起があり、左右の耳に1本ずつ対応している。2つの入力を別々の枝に分けることで、ニューロンは「AND」をより厳しくする。1本の枝に2つの電流が入ると互いに邪魔し合うからだ。

運動学習回路の小さなニューロンには、短い樹状突起が約4本あり、それぞれが自分の入力をつかむ。こうしたニューロンは数百億個あり、それぞれがたった4つの信号を足すだけだ。その代わり全体として入力を数百万通りの組み合わせに展開し、10万の入力をもつ大きな出力ニューロンがその中から必要なものを選ぶ。

最後に、巨大な木と数千の入力をもつニューロンがある。加算器であり分類器であり、個々の枝はほとんど独立した小さな回路のように働く。

\section*{なぜそうなるのか}

すべてを決めるのは充電である。入力の少ない小さなニューロンは、1つの大きな接続部から1ミリ秒未満で充電される。こうしたニューロンは時間的に正確で、信号を確実に先へ伝える。大きな木は、1つの接続部では決して充電されない。その前に電荷が漏れてしまう。数十の信号が一緒に届く必要があるが、その代わりそれらを足し合わせ、比べることができる。入力が少なく接続部が大きいのは正確なタイミングのため。入力が多いのは数えるためである。

\pencilfig{128}{%
% charging: a small neuron reaches threshold from one large synapse at once; a big tree barely moves
% from one input and needs many inputs arriving together
\foreach \dx/\t in {0/{入力が少ない}, 4.3/{大きな木}}{
  \draw[pen] (\dx,0) -- (\dx+3.6,0);
  \draw[pcGray, densely dashed, line width=0.5pt] (\dx,0.9) -- (\dx+3.6,0.9);
  \node[font=\small\itshape, text=pcGray, anchor=south] at (\dx+1.8,1.75) {\t};}
\node[lbl, anchor=east] at (-0.05,0.9) {閾値};
% small neuron
\draw[pcBlue, line width=2pt] (0.6,-0.15) -- (0.6,-0.6);
\draw[penlite=pcBlue] (0,0) -- (0.6,0) -- plot[domain=0.6:0.8, samples=12] (\x,{1.3*(1-exp(-(\x-0.6)/0.18))}) -- (0.8,1.6) -- (0.84,0) -- (3.5,0);
\node[lbl, text=pcRed, anchor=west] at (0.85,1.45) {すぐにカチッ};
\node[lbl, text=pcBlue, anchor=north] at (0.6,-0.62) {大きな入力1つ};
% big tree
\draw[pcBlue, line width=0.6pt] (4.8,-0.15) -- (4.8,-0.45);
\foreach \s in {6.15,6.2,6.25,6.3,6.35,6.4,6.45,6.5}{\draw[pcBlue, line width=0.6pt] (\s,-0.15) -- (\s,-0.45);}
\draw[penlite=pcBlue] (4.3,0) -- (4.8,0) -- plot[domain=4.8:6.15, samples=20] (\x,{0.16*((\x-4.8)/0.15)*exp(1-(\x-4.8)/0.15)}) -- plot[domain=6.15:6.62, samples=14] (\x,{1.25*(1-exp(-(\x-6.15)/0.4))}) -- (6.62,1.6) -- (6.66,0) -- (7.9,0);
\node[lbl, anchor=south] at (4.95,0.2) {ほぼ変化なし};
\node[lbl, text=pcBlue, anchor=north] at (6.33,-0.47) {多数が同時に};
}{小さなニューロンは1つの大きな接続部から1ミリ秒未満で充電され、すぐにカチッと鳴る。大きな木は入力1つではほとんど変わらない。数十の信号が一緒に届く必要がある。}

こうしてニューロンの3つ目の仕分けが得られた。物質による、素子による、そして入力の数による仕分けである。それぞれが同じマシンを新しい側面から見ている。

\fullversion{19}
